\chapter{Implémentation / Réalisation}
\label{chap:implementation}

\section{Environnement de travail}

\subsection{Environnement logiciel}

\begin{itemize}
  \item \textbf{Langage backend~:} Python 3.10
  \item \textbf{Framework web~:} FastAPI \cite{fastapi_docs}, serveur
        ASGI Uvicorn
  \item \textbf{Base de données applicative~:} MySQL 8
  \item \textbf{Client Wazuh Indexer~:} \texttt{opensearch-py}
  \item \textbf{Frontend~:} Vue 3 \cite{vuejs_docs} (Composition API),
        Pinia (gestion d'état), Vue Router, Tailwind CSS, Vite
        (\textit{build tool})
  \item \textbf{Fournisseurs de LLM~:} toute API compatible
        \texttt{/chat/completions} au format OpenAI
        \cite{openai_toolcalling} (Groq utilisé en développement), ou un
        modèle auto-hébergé via Ollama \cite{ollama_docs}
  \item \textbf{Enrichissement externe~:} API AbuseIPDB \cite{abuseipdb}
        pour la réputation d'adresses IP
  \item \textbf{Gestion de versions~:} Git
\end{itemize}

\subsection{Environnement matériel}

Le développement et les tests ont été réalisés sur une machine de
développement standard (architecture x86-64, système Linux), le Wazuh
Indexer et la base MySQL étant déployés localement pour les besoins du
développement. Le caractère \textit{stateless} de l'API et l'absence de
traitement lourd côté serveur (le calcul est délégué au fournisseur de LLM)
rendent le déploiement peu exigeant en ressources~; un serveur disposant de
quelques gigaoctets de mémoire suffit pour l'exécution de l'API, de la base
de données et, le cas échéant, d'un modèle local de taille modeste via
Ollama.

\section{Description et traitement des données}

Deux sources de données coexistent dans la plateforme~:

\begin{description}
  \item[Les alertes Wazuh] sont des documents JSON stockés dans le Wazuh
    Indexer, structurés notamment autour des champs \texttt{rule} (identifiant
    et niveau de sévérité de la règle déclenchée), \texttt{agent}
    (identifiant et nom de l'hôte concerné), \texttt{data} (champs extraits
    par les décodeurs Wazuh, tels que \texttt{srcip}/\texttt{dstip}) et
    \texttt{full\_log} (la ligne de journal brute). Ces documents sont
    interrogés en lecture seule~; la plateforme ne modifie jamais les
    alertes elles-mêmes.
  \item[Les données applicatives] sont stockées dans une base MySQL
    relationnelle (schéma détaillé en \cref{fig:erd}), couvrant les
    investigations, les preuves collectées, les analyses produites, les
    actions de l'IA, les actifs, les politiques, les utilisateurs et les
    conversations de l'assistant. Les tables sont créées et migrées
    automatiquement au démarrage du serveur.
\end{description}

\section{Modules principaux}

\subsection{Collecte et politiques (\texttt{poller\_service.py},
\texttt{investigation\_policy.py})}

Le \textbf{Poller Service} interroge le Wazuh Indexer à intervalle
configurable, en ne retenant que les alertes postérieures au dernier
horodatage traité (mécanisme de \textit{watermark} persistant, pour éviter
les doublons ou les pertes entre deux cycles). Chaque alerte récupérée est
soumise au \textbf{Policy Engine}, qui évalue un ensemble de règles
configurables (sévérité minimale, groupes de règles, agents concernés) afin
de décider si une investigation doit être ouverte.

\subsection{Gestion des investigations (\texttt{investigation\_manager.py})}

L'\textbf{Investigation Manager} crée l'enregistrement d'investigation,
recherche ou crée automatiquement l'\textbf{actif} correspondant à l'agent
Wazuh concerné, puis transmet l'investigation à l'AI Investigator.

\subsection{Agent d'investigation (\texttt{ai\_investigator/investigator.py})}

C'est le cœur du projet. Étant donné une alerte déclenchante, l'agent~:

\begin{enumerate}
  \item récupère un \textbf{contexte de référence} (fiche de l'actif
        concerné et historique récent d'alertes du même agent) avant même
        de commencer à raisonner, pour répondre immédiatement à la question
        \og ce comportement est-il habituel pour cet hôte~?\fg{} sans appel
        d'outil supplémentaire~;
  \item engage une \textbf{boucle agentique} (inspirée du paradigme ReAct
        \cite{react2022}, \cref{chap:etat_de_lart})~: à chaque étape, le
        modèle de langage reçoit l'historique de la conversation et la
        liste des outils disponibles, et peut soit appeler un outil, soit
        rendre son verdict~;
  \item lorsqu'un outil est appelé, celui-ci est exécuté, son résultat est
        journalisé (table \texttt{investigation\_actions}, avec le
        raisonnement du modèle au moment de l'appel) et réinjecté dans la
        conversation~;
  \item la boucle se termine lorsque le modèle produit un verdict, ou
        lorsqu'un budget d'étapes est atteint, auquel cas un verdict final
        est forcé (\texttt{tool\_choice="none"}) sur la base des preuves
        déjà réunies.
\end{enumerate}

\paragraph{Garde-fous de la boucle agentique.} Un modèle de langage
n'étant pas un composant strictement déterministe ni intrinsèquement borné
en coût, plusieurs mécanismes de protection ont été mis en place~:

\begin{itemize}
  \item un nombre maximal d'itérations (\texttt{MAX\_STEPS})~;
  \item une taille maximale par résultat d'outil et pour la conversation
        totale renvoyée au modèle (\texttt{MAX\_TOOL\_RESULT\_CHARS},
        \texttt{MAX\_CONVERSATION\_CHARS}), les appels \textit{chat
        completions} étant \textit{stateless} et donc facturés/limités sur
        l'intégralité du contexte renvoyé à chaque tour~;
  \item une \textbf{détection des appels d'outils dupliqués}~: un appel
        strictement identique (même outil, mêmes arguments) à un appel déjà
        effectué dans la même investigation n'est pas réexécuté, son
        résultat déjà connu est simplement rappelé au modèle~;
  \item une \textbf{normalisation des arguments} reçus du modèle~: un LLM
        peut renvoyer un entier sous forme de chaîne de caractères
        (\texttt{"50"} au lieu de \texttt{50}) dans son appel de fonction ;
        chaque outil convertit automatiquement ses arguments selon son
        propre schéma JSON avant exécution (\cref{lst:coerce})~;
  \item des \textbf{délais d'expiration réseau} explicites sur les appels
        au fournisseur de LLM et au Wazuh Indexer~;
  \item une exécution de l'investigation entièrement encadrée par une
        gestion d'erreurs~: toute exception (LLM, outil, base de données)
        aboutit à un verdict \texttt{ERROR} enregistré plutôt qu'à une
        investigation bloquée indéfiniment à l'état \og en cours\fg.
\end{itemize}

\begin{lstlisting}[language=Python, caption={Normalisation générique des
arguments d'un outil selon son schéma JSON}, label={lst:coerce}]
def coerce_arguments(self, arguments: dict) -> dict:
    """Cast argument values to match this tool's declared JSON
    Schema types, e.g. {"limit": "50"} -> {"limit": 50}."""
    props = self.parameters.get("properties", {})
    coerced = dict(arguments)
    for key, spec in props.items():
        if key not in coerced or not isinstance(coerced[key], str):
            continue
        expected = spec.get("type")
        raw = coerced[key].strip()
        try:
            if expected == "integer":
                coerced[key] = int(raw)
            elif expected == "number":
                coerced[key] = float(raw)
        except ValueError:
            pass  # left as-is; execute() surfaces the problem
    return coerced
\end{lstlisting}

\subsection{Outils disponibles pour l'IA}

Le \cref{tab:tools} liste les outils partagés par l'agent d'investigation
et l'assistant conversationnel. Ils sont volontairement en \textbf{lecture
seule}.

\begin{table}[H]
\centering
\caption{Outils exposés à l'IA}
\label{tab:tools}
\begin{tabular}{@{}p{3.6cm}p{3cm}p{7cm}@{}}
\toprule
\textbf{Outil} & \textbf{Source} & \textbf{Rôle} \\
\midrule
\texttt{search\_alerts} & Wazuh Indexer & Rechercher des alertes par
agent, IP, règle, plage temporelle, texte libre \\
\texttt{count\_alerts} & Wazuh Indexer & Compter exactement les alertes
correspondant à des filtres, sans les limites de \texttt{search\_alerts} \\
\texttt{get\_alert\_details} & Wazuh Indexer & Détail complet d'une
alerte \\
\texttt{search\_investigations} & MySQL & Rechercher des investigations
passées \\
\texttt{count\_investigations} & MySQL & Compter exactement les
investigations correspondant à des filtres \\
\texttt{get\_investigation}, \texttt{get\_investigation\_evidence},
\texttt{get\_investigation\_analysis} & MySQL & Détail, preuves et
analyses d'une investigation \\
\texttt{search\_assets}, \texttt{get\_asset} & MySQL & Rechercher/consulter
un actif (hôte) \\
\texttt{check\_ip\_reputation} & AbuseIPDB \cite{abuseipdb} & Score de
réputation d'une adresse IP publique (score de confiance d'abus, nombre de
signalements, pays, fournisseur) \\
\bottomrule
\end{tabular}
\end{table}

\subsection{Assistant conversationnel (\texttt{ai\_investigator/chatbot.py})}

L'assistant conversationnel réutilise exactement le même ensemble d'outils
que l'agent d'investigation (\cref{fig:classdiag}), mais dans un cadre
multi-tours en libre échange plutôt que centré sur une alerte unique. Pour
éviter que la taille de la conversation — et donc le coût des appels au
LLM — ne croisse indéfiniment au fil d'une session de discussion, seule la
question de l'utilisateur et la réponse finale de l'assistant sont
persistées~: l'éventuel aller-retour d'appels d'outils au sein d'un tour de
conversation est reconstruit à partir de l'historique, puis abandonné une
fois la réponse produite.

\subsection{Authentification (\texttt{api/auth.py})}

Conformément à la conception présentée en \cref{sec:securite_conception},
la vérification des identifiants (\cref{lst:auth}) consiste en une simple
requête HTTP authentifiée vers le Wazuh Indexer~: une réponse \texttt{200}
confirme la validité du couple identifiant/mot de passe, une réponse
\texttt{401}/\texttt{403} l'infirme. Le mot de passe n'est jamais persisté.

\begin{lstlisting}[language=Python, caption={Vérification des identifiants
auprès du Wazuh Indexer}, label={lst:auth}]
def verify_wazuh_credentials(username, password, *,
                              host, port, verify_ssl,
                              timeout=5.0) -> bool:
    try:
        resp = requests.get(
            f"https://{host}:{port}/",
            auth=(username, password),
            verify=verify_ssl,
            timeout=timeout,
        )
    except requests.exceptions.RequestException:
        return False
    if resp.status_code == 200:
        return True
    if resp.status_code in (401, 403):
        return False
    return False  # erreur cote Indexer, pas une preuve d'invalidite
\end{lstlisting}

Un jeton JWT signé (durée de vie configurable, 12 heures par défaut) est
ensuite émis et doit accompagner chaque requête suivante
(\textit{header} \texttt{Authorization: Bearer <token>}). Toutes les routes
de l'API sont protégées par ce mécanisme, à l'exception de la route de
connexion elle-même et du point de contrôle de santé (\textit{health
check}).

\section{Interface utilisateur}

L'interface web (Vue 3) propose les écrans suivants~: un \textbf{tableau de
bord} synthétisant l'état des investigations en cours, une \textbf{liste
des investigations} filtrable, une \textbf{vue détail} par investigation
(chronologie des actions de l'IA, preuves, verdict), une \textbf{gestion
des politiques} et des \textbf{actifs}, un \textbf{assistant conversationnel}
(historique de sessions, affichage du raisonnement et de la trace des
outils appelés lorsqu'ils sont disponibles), et un \textbf{écran de
connexion}.

\todofig{Capture d'écran -- Tableau de bord}
\todofig{Capture d'écran -- Détail d'une investigation (chronologie des
actions de l'IA et verdict)}
\todofig{Capture d'écran -- Assistant conversationnel}
